%! Author = sriadityavedantam
%! Date = 3/30/26

\section{Chapter 2}

\begin{problem}{
    Suppose $\gcd(a,b)=1$.
    If $a\mid c$ and $b\mid c$, prove that $ab\mid c$.
}
    Since
    \[
        \gcd(a,b)=1,
    \]
    there exist $x,y\in\z$ such that
    \[
        1=ax+by.
    \]
    Now multiply both sides by $c$.
    Then
    \[
        c=acx+bcy.
    \]
    Since $a\mid c$ and $b\mid c$, there exist $m,n\in\z$ such that
    \[
        c=am
        \qquad\text{and}\qquad
        c=bn.
    \]
    Substitute these into the previous equation:
    \begin{align*}
        c &= a(bn)x+b(am)y\\
        &= ab(nx+my).
    \end{align*}
    Let
    \[
        t\coloneqq nx+my\in\z.
    \]
    Then
    \[
        c=abt.
    \]
    Hence
    \[
        ab\mid c.
    \]
\end{problem}

\begin{problem}{
    If $\gcd(a,c)=1$ and $\gcd(b,c)=1$, prove that $\gcd(ab,c)=1$.
}
    Since
    \[
        \gcd(a,c)=1,
    \]
    there exist $x,y\in\z$ such that
    \[
        1=ax+cy.
    \]
    Since
    \[
        \gcd(b,c)=1,
    \]
    there exist $m,n\in\z$ such that
    \[
        1=bm+cn.
    \]
    Multiply the first equation by $b$:
    \[
        b=abx+bcy.
    \]
    Now multiply the second equation by $cy$:
    \[
        cy=b(cmy)+c^{2}ny.
    \]
    Substitute the expression
    \[
        b=abx+bcy
    \]
    into
    \[
        1=bm+cn:
    \]
    \begin{align*}
        1 &= (abx+bcy)m+cn\\
        &= ab(mx)+c(bym+n).
    \end{align*}
    Let
    \[
        p\coloneqq mx
        \qquad\text{and}\qquad
        q\coloneqq bym+n.
    \]
    Then
    \[
        1=abp+cq.
    \]
    Therefore
    \[
        \gcd(ab,c)=1.
    \]
\end{problem}

\begin{problem}{
    (a) If $a,b,u,v\in\z$ are such that $au+bv=1$, prove that $\gcd(a,b)=1$.

    (b) Show by example that if $au+bv=d>1$, then $\gcd(a,b)$ may not be $d$.
}
    \textbf{(a).}
    Assume
    \[
        au+bv=1.
    \]
    Let $d=\gcd(a,b)$.
    Then
    \[
        d\mid a
        \qquad\text{and}\qquad
        d\mid b.
    \]
    So $d$ divides every linear combination of $a$ and $b$.
    In particular,
    \[
        d\mid (au+bv)=1.
    \]
    Hence
    \[
        d=1.
    \]
    Therefore
    \[
        \gcd(a,b)=1.
    \]

    \textbf{(b).}
    Take
    \[
        a=6,
        \qquad
        b=9,
        \qquad
        u=2,
        \qquad
        v=-1.
    \]
    Then
    \[
        au+bv=6(2)+9(-1)=12-9=3.
    \]
    So here
    \[
        au+bv=d=3>1.
    \]
    But
    \[
        \gcd(6,9)=3,
    \]
    so this still equals $d$.
    To get an example where the gcd is not $d$, take instead
    \[
        a=6,
        \qquad
        b=9,
        \qquad
        u=1,
        \qquad
        v=1.
    \]
    Then
    \[
        au+bv=6+9=15.
    \]
    So
    \[
        d=15>1,
    \]
    but
    \[
        \gcd(6,9)=3\neq 15.
    \]
    Thus if
    \[
        au+bv=d>1,
    \]
    it need not follow that
    \[
        \gcd(a,b)=d.
    \]
\end{problem}

\begin{problem}{
    If $a\mid c$ and $b\mid c$ and $\gcd(a,b)=d$, prove that $ab\mid cd$.
}
    Let
    \[
        \gcd(a,b)=d.
    \]
    Then there exist $a_1,b_1\in\z$ such that
    \[
        a=da_1
        \qquad\text{and}\qquad
        b=db_1,
    \]
    with
    \[
        \gcd(a_1,b_1)=1.
    \]
    Since
    \[
        a\mid c
        \qquad\text{and}\qquad
        b\mid c,
    \]
    we have
    \[
        da_1\mid c
        \qquad\text{and}\qquad
        db_1\mid c.
    \]
    Hence
    \[
        a_1\mid \frac{c}{d}
        \qquad\text{and}\qquad
        b_1\mid \frac{c}{d}.
    \]
    Since
    \[
        \gcd(a_1,b_1)=1,
    \]
    the previous problem implies
    \[
        a_{1}b_1\mid \frac{c}{d}.
    \]
    So there exists $t\in\z$ such that
    \[
        \frac{c}{d}=a_{1}b_{1}t.
    \]
    Multiply both sides by $d^2$:
    \[
        cd=d^{2}a_{1}b_{1}t=(da_1)(db_1)t=abt.
    \]
    Therefore
    \[
        ab\mid cd.
    \]
\end{problem}

\begin{problem}{
    If $a>0$ and $b>0$, prove that
    \[
        lcm[a,b]=\frac{ab}{\gcd(a,b)}.
    \]
}
    Let
    \[
        d=\gcd(a,b).
    \]
    Then there exist $x,y\in\z_{>0}$ such that
    \[
        a=dx,
        \qquad
        b=dy,
    \]
    and
    \[
        \gcd(x,y)=1.
    \]
    We claim that
    \[
        lcm[a,b]=dxy.
    \]
    First, $dxy$ is a common multiple of $a$ and $b$ because
    \[
        dxy=ay=bx.
    \]
    So
    \[
        a\mid dxy
        \qquad\text{and}\qquad
        b\mid dxy.
    \]
    Now let $m$ be any common multiple of $a$ and $b$.
    Then
    \[
        a\mid m
        \qquad\text{and}\qquad
        b\mid m.
    \]
    So
    \[
        dx\mid m
        \qquad\text{and}\qquad
        dy\mid m.
    \]
    Hence
    \[
        x\mid \frac{m}{d}
        \qquad\text{and}\qquad
        y\mid \frac{m}{d}.
    \]
    Since
    \[
        \gcd(x,y)=1,
    \]
    it follows that
    \[
        xy\mid \frac{m}{d}.
    \]
    Thus
    \[
        dxy\mid m.
    \]
    So $dxy$ divides every common multiple of $a$ and $b$, which means it is the least common multiple.
    Therefore
    \[
        lcm[a,b]=dxy.
    \]
    But
    \[
        \frac{ab}{\gcd(a,b)}
        =
        \frac{(dx)(dy)}{d}
        =
        dxy.
    \]
    Hence
    \[
        lcm[a,b]=\frac{ab}{\gcd(a,b)}.
    \]
\end{problem}

\begin{problem}{
    Prove that:

    (a)
    \[
        \gcd(a,b)\mid \gcd(a+b,a-b).
    \]

    (b)
    If $a$ is odd and $b$ is even, then
    \[
        \gcd(a,b)=\gcd(a+b,a-b).
    \]

    (c)
    If $a$ and $b$ are odd, then
    \[
        2\,\gcd(a,b)=\gcd(a+b,a-b).
    \]
}
    \textbf{(a).}
    Let
    \[
        d=\gcd(a,b).
    \]
    Then
    \[
        d\mid a
        \qquad\text{and}\qquad
        d\mid b.
    \]
    Hence
    \[
        d\mid (a+b)
        \qquad\text{and}\qquad
        d\mid (a-b).
    \]
    Therefore $d$ is a common divisor of $a+b$ and $a-b$, so
    \[
        d\mid \gcd(a+b,a-b).
    \]

    \textbf{(b).}
    Let
    \[
        d=\gcd(a,b).
    \]
    By part (a),
    \[
        d\mid \gcd(a+b,a-b).
    \]
    Now let
    \[
        e=\gcd(a+b,a-b).
    \]
    Then
    \[
        e\mid (a+b)
        \qquad\text{and}\qquad
        e\mid (a-b).
    \]
    So
    \[
        e\mid \big((a+b)+(a-b)\big)=2a
    \]
    and
    \[
        e\mid \big((a+b)-(a-b)\big)=2b.
    \]
    Since $a$ is odd and $b$ is even, any common divisor of $a+b$ and $a-b$ must be odd.
    Thus
    \[
        \gcd(e,2)=1.
    \]
    Because
    \[
        e\mid 2a
    \]
    and
    \[
        \gcd(e,2)=1,
    \]
    it follows that
    \[
        e\mid a.
    \]
    Similarly,
    \[
        e\mid b.
    \]
    Hence
    \[
        e\mid \gcd(a,b)=d.
    \]
    Together with part (a), this gives
    \[
        \gcd(a,b)=\gcd(a+b,a-b).
    \]

    \textbf{(c).}
    Let
    \[
        d=\gcd(a,b).
    \]
    Since $a$ and $b$ are odd, write
    \[
        a=da_1
        \qquad\text{and}\qquad
        b=db_1,
    \]
    where
    \[
        \gcd(a_1,b_1)=1.
    \]
    Because $a$ and $b$ are odd, both $a_1$ and $b_1$ are odd.
    Then
    \[
        a+b=d(a_1+b_1)
        \qquad\text{and}\qquad
        a-b=d(a_1-b_1).
    \]
    So
    \[
        \gcd(a+b,a-b)=d\cdot \gcd(a_1+b_1,a_1-b_1).
    \]
    Now $a_1+b_1$ and $a_1-b_1$ are both even, so
    \[
        2\mid \gcd(a_1+b_1,a_1-b_1).
    \]
    Also, if an odd integer divides both $a_1+b_1$ and $a_1-b_1$, then it divides their sum and difference:
    \begin{gather*}
        (a_1+b_1)+(a_1-b_1)=2a_1,\\
    (a_1+b_1)-(a_1-b_1)=2b_1.\\
    \end{gather*}
    Since the divisor is odd, it must divide both $a_1$ and $b_1$.
    But
    \[
        \gcd(a_1,b_1)=1,
    \]
    so no odd divisor greater than $1$ can occur.
    Hence
    \[
        \gcd(a_1+b_1,a_1-b_1)=2.
    \]
    Therefore
    \[
        \gcd(a+b,a-b)=2d=2\,\gcd(a,b).
    \]
\end{problem}

\begin{problem}{
    Use the GCD algorithm to show that the GCD of $166$ and $39$ is $1$ and to find integers $x$ and $y$ such that
    \[
        1=x166+y39.
    \]
    Show your work.
}
    Apply the Euclidean algorithm:
    \begin{align*}
        166 &= 4\cdot 39 + 10,\\
        39 &= 3\cdot 10 + 9,\\
        10 &= 1\cdot 9 + 1,\\
        9 &= 9\cdot 1 + 0.
    \end{align*}
    Hence
    \[
        \gcd(166,39)=1.
    \]
    Now work backward to find the extended gcd:
    \begin{align*}
        1 &= 10-9,\\
        &= 10-(39-3\cdot 10),\\
        &= 4\cdot 10-39,\\
        &= 4(166-4\cdot 39)-39,\\
        &= 4\cdot 166-17\cdot 39.
    \end{align*}
    Thus
    \[
        1=4\cdot 166+(-17)\cdot 39.
    \]
    So
    \[
        x=4
        \qquad\text{and}\qquad
        y=-17.
    \]
\end{problem}

\begin{problem}{
    Let $p$ be an integer other than $0,\pm 1$.
    Prove that $p$ is prime if and only if for each $a\in\z$ either $\gcd(a,p)=1$ or $p\mid a$.
}
    \textbf{($\implies$).}
    Assume that $p$ is prime.
    Let $a\in\z$.
    If $p\mid a$, then we are done.
    So suppose $p\nmid a$.
    We must show that
    \[
        \gcd(a,p)=1.
    \]
    Let
    \[
        d=\gcd(a,p).
    \]
    Then
    \[
        d\mid p.
    \]
    Since $p$ is prime, its only divisors are
    \[
        \pm 1,\pm p.
    \]
    Because $p\nmid a$, we cannot have $d=\pm p$.
    Thus
    \[
        d=\pm 1.
    \]
    Since the gcd is positive, it follows that
    \[
        \gcd(a,p)=1.
    \]

    \textbf{($\impliedby$).}
    Assume that for each $a\in\z$, either
    \[
        \gcd(a,p)=1
        \quad\text{or}\quad
        p\mid a.
    \]
    We prove that $p$ is prime.
    Let $d\in\z$ with
    \[
        d\mid p.
    \]
    Then there exists $m\in\z$ such that
    \[
        p=dm.
    \]
    Apply the given property to $a=d$.
    Either
    \[
        \gcd(d,p)=1
        \quad\text{or}\quad
        p\mid d.
    \]
    If
    \[
        \gcd(d,p)=1,
    \]
    then since $d\mid p$, the only positive possibility is
    \[
        d=1.
    \]
    If
    \[
        p\mid d,
    \]
    then there exists $t\in\z$ such that
    \[
        d=pt.
    \]
    Since also
    \[
        d\mid p,
    \]
    we must have
    \[
        d=\pm p.
    \]
    Thus every divisor of $p$ is one of
    \[
        \pm 1,\pm p.
    \]
    Since $p\neq 0,\pm 1$, it follows that $p$ is prime.
\end{problem}

\begin{problem}{
    Let $p$ be an integer other than $0,\pm 1$ with this property:
    Whenever $b$ and $c$ are integers such that $p\mid bc$, then $p\mid b$ or $p\mid c$.
    Prove that $p$ is prime.
}
    Suppose $d\mid p$.
    Then there exists $t\in\z$ such that
    \[
        p=dt.
    \]
    Since
    \[
        p\mid dt,
    \]
    the given property implies that
    \[
        p\mid d
        \quad\text{or}\quad
        p\mid t.
    \]
    If
    \[
        p\mid d,
    \]
    then
    \[
        d=\pm p.
    \]
    If
    \[
        p\mid t,
    \]
    then
    \[
        t=\pm p
    \]
    or at least $|t|\geq |p|$, and since
    \[
        p=dt,
    \]
    this forces
    \[
        d=\pm 1.
    \]
    So every divisor of $p$ is one of
    \[
        \pm 1,\pm p.
    \]
    Since $p\neq 0,\pm 1$, it follows that $p$ is prime.
\end{problem}

\begin{problem}{
    If
    \[
        a=p_1^{r_1}p_2^{r_2}\dots p_k^{r_k}
    \]
    and
    \[
        b=p_1^{s_1}p_2^{s_2}\dots p_k^{s_k},
    \]
    where $p_1,p_2,\dots,p_k$ are distinct positive primes and each $r_i,s_i\geq 0$, then prove that:

    \textbf{(a)}
    \[
        \gcd(a,b)=p_1^{n_1}p_2^{n_2}\dots p_k^{n_k},
    \]
    where for each $i$,
    \[
        n_i=\min\{r_i,s_i\}.
    \]

    \textbf{(b)}
    \[
        lcm[a,b]=p_1^{t_1}p_2^{t_2}\dots p_k^{t_k},
    \]
    where for each $i$,
    \[
        t_i=\max\{r_i,s_i\}.
    \]
}
    \textbf{(a).}
    Let
    \[
        d\coloneqq p_1^{n_1}p_2^{n_2}\dots p_k^{n_k},
    \]
    where
    \[
        n_i=\min\{r_i,s_i\}.
    \]
    We show that $d=\gcd(a,b)$.

    First, $d\mid a$ and $d\mid b$.
    Indeed, for each $i$,
    \[
        n_i\leq r_i
        \quad\text{and}\quad
        n_i\leq s_i,
    \]
    so every prime power appearing in $d$ appears in both $a$ and $b$ with exponent at least $n_i$.
    Now let $q$ be any common divisor of $a$ and $b$.
    Write
    \[
        q=p_1^{v_1}p_2^{v_2}\dots p_k^{v_k},
    \]
    where each $v_i\geq 0$.
    Since $q\mid a$, we must have
    \[
        v_i\leq r_i
    \]
    for each $i$.
    Since $q\mid b$, we must also have
    \[
        v_i\leq s_i
    \]
    for each $i$.
    Hence
    \[
        v_i\leq \min\{r_i,s_i\}=n_i
    \]
    for every $i$.
    Therefore
    \[
        q\mid d.
    \]
    So $d$ is a common divisor of $a$ and $b$, and every common divisor divides $d$.
    Hence
    \[
        \gcd(a,b)=d=p_1^{n_1}p_2^{n_2}\dots p_k^{n_k}.
    \]

    \textbf{(b).}
    Let
    \[
        t_i=\max\{r_i,s_i\}
    \]
    for each $i$, and define
    \[
        m\coloneqq p_1^{t_1}p_2^{t_2}\dots p_k^{t_k}.
    \]
    We show that
    \[
        lcm[a,b]=m.
    \]
    First, $a\mid m$ and $b\mid m$.
    Indeed, for each $i$,
    \[
        r_i\leq t_i
        \quad\text{and}\quad
        s_i\leq t_i.
    \]
    So $m$ is a common multiple of $a$ and $b$.

    Now let $N$ be any common multiple of $a$ and $b$.
    Write
    \[
        N=p_1^{w_1}p_2^{w_2}\dots p_k^{w_k}\cdot u,
    \]
    where $u$ is not divisible by any of the primes $p_i$.
    Since $a\mid N$, we have
    \[
        r_i\leq w_i
    \]
    for each $i$.
    Since $b\mid N$, we also have
    \[
        s_i\leq w_i
    \]
    for each $i$.
    Therefore
    \[
        t_i=\max\{r_i,s_i\}\leq w_i
    \]
    for every $i$.
    So
    \[
        m\mid N.
    \]

    Thus $m$ is a common multiple of $a$ and $b$, and it divides every common multiple.
    Hence
    \[
        lcm[a,b]=m=p_1^{t_1}p_2^{t_2}\dots p_k^{t_k}.
    \]
\end{problem}

\begin{problem}{
    Prove that
    \[
        a\mid b
        \quad\text{if and only if}\quad
        a^2\mid b^2.
    \]
}
    \textbf{($\implies$).}
    Suppose
    \[
        a\mid b.
    \]
    Then there exists $x\in\z$ such that
    \[
        b=ax.
    \]
    Squaring both sides gives
    \[
        b^2=a^{2}x^2.
    \]
    Hence
    \[
        a^2\mid b^2.
    \]
    \textbf{($\impliedby$).}
    Suppose
    \[
        a^2\mid b^2.
    \]
    We want to show that
    \[
        a\mid b.
    \]
    Write the prime factorizations of $a$ and $b$.
    If
    \[
        a=\pm p_1^{r_1}\cdots p_k^{r_k},
    \]
    then
    \[
        a^2=p_1^{2r_1}\cdots p_k^{2r_k}.
    \]
    If
    \[
        a^2\mid b^2,
    \]
    then for each prime $p_i$, the exponent of $p_i$ in $b^2$ is at least $2r_i$.
    So the exponent of $p_i$ in $b$ is at least $r_i$.
    Therefore every prime power dividing $a$ also divides $b$, and hence
    \[
        a\mid b.
    \]
    Thus
    \[
        a\mid b
        \quad\text{if and only if}\quad
        a^2\mid b^2.
    \]
\end{problem}

\begin{problem}{
    Let $p$ be prime and $1<k<p$.
    Prove that $p$ divides the binomial coefficient
    \[
        \binom{p}{k}.
    \]
}
    Recall that
    \[
        \binom{p}{k}=\frac{p!}{k!(p-k)!}.
    \]
    Write
    \[
        p!=p(p-1)!.
    \]
    Then
    \[
        \binom{p}{k}=p\cdot \frac{(p-1)!}{k!(p-k)!}.
    \]
    We now show that
    \[
        \frac{(p-1)!}{k!(p-k)!}\in\z.
    \]
    But this is automatic, since
    \[
        \binom{p}{k}
    \]
    is an integer.

    Also, because
    \[
        1<k<p,
    \]
    both $k!$ and $(p-k)!$ are products of positive integers all strictly less than $p$.
    Since $p$ is prime, none of those factors is divisible by $p$.
    So the factor of $p$ in the numerator does not cancel.

    Therefore
    \[
        p\mid \binom{p}{k}.
    \]
\end{problem}

\begin{problem}{
    Let $p,q$ be primes with $p\geq 5$ and $q\geq 5$.
    Prove that
    \[
        24\mid (p^2-q^2).
    \]
}
    Since $p,q\geq 5$ are prime, both are odd.
    Thus
    \[
        p^2-q^2=(p-q)(p+q).
    \]
    We first show that
    \[
        8\mid (p^2-q^2).
    \]
    Every odd integer is congruent to either $1$, $3$, $5$, or $7$ modulo $8$, and in each case its square is congruent to $1$ modulo $8$.
    So
    \[
        p^2\equiv 1\pmod{8}
        \quad\text{and}\quad
        q^2\equiv 1\pmod{8}.
    \]
    Hence
    \[
        p^2-q^2\equiv 0\pmod{8}.
    \]
    Therefore
    \[
        8\mid (p^2-q^2).
    \]

    Next we show that
    \[
        3\mid (p^2-q^2).
    \]
    Since $p$ and $q$ are primes at least $5$, neither is divisible by $3$.
    So each is congruent to either $1$ or $-1$ modulo $3$.
    Hence
    \[
        p^2\equiv 1\pmod{3}
        \quad\text{and}\quad
        q^2\equiv 1\pmod{3}.
    \]
    Therefore
    \[
        p^2-q^2\equiv 0\pmod{3},
    \]
    so
    \[
        3\mid (p^2-q^2).
    \]
    Since
    \[
        \gcd(8,3)=1,
    \]
    it follows that
    \[
        24=8\cdot 3 \mid (p^2-q^2).
    \]
    Hence
    \[
        24\mid (p^2-q^2).
    \]
\end{problem}